\section{Проектирование механизма миграции}\label{sec:design}

\subsection{Архитектура миграции с даунтаймом}\label{subsec:mvp-architecture}

Архитектура построена вокруг бакета и его статусной модели в \textit{service-staas} (рис.~\ref{fig:bucket-states}).
Помимо обычных состояний (\code{READY}, \code{DELETING} и т.\,п.) для бакета вводится семейство статусов \code{MIGRATION\_*}, отражающих этапы переноса.
По принадлежности бакета к этому семейству остальные части платформы определяют, что идёт миграция, и подстраивают своё поведение~--- например, блокируют конкурирующие операции или иначе обрабатывают запросы.

\begin{figure}[H]
    \centering
    \includegraphics[width=0.7\textwidth]{bucket-states}
    \caption{Статусная модель бакета в \textit{service-staas}}
    \label{fig:bucket-states}
\end{figure}

Архитектура этого этапа складывается из оркестратора, ведущего бакет по статусной модели, и сменного движка копирования.

\subsubsection{Оркестратор}\label{subsubsec:orchestrator}

Оркестратор~--- это логический компонент \textit{service-staas}, отвечающий за продвижение бакета по статусной модели.
Он забирает бакеты в активных миграционных статусах и проводит каждый по фазам, изображённым на рис.~\ref{fig:bucket-states}:
\begin{itemize}
    \item \code{MIGRATION\_SCHEDULED}~--- создание целевого бакета и выдача доступов;
    \item \code{MIGRATION\_IN\_PROGRESS}~--- копирование данных;
    \item \code{MIGRATION\_VERIFYING}~--- сверка содержимого исходного и целевого бакетов (множества ключей и время последней модификации);
    \item \code{MIGRATION\_FINISHED}~--- переключение DNS-записи в \gls{consul} и постановка исходного бакета на удаление;
    \item \code{MIGRATION\_FAILED}~--- терминальное состояние, выставляется автоматически, если на фазе \code{MIGRATION\_VERIFYING} сверка не прошла успешно; повторный запуск инициируется ручным возвратом в \code{MIGRATION\_SCHEDULED}.
\end{itemize}

Остановка записи остаётся за пределами автоматики: она согласуется с владельцем бакета через внешний канал (тикет, чат).

Единственным внешним компонентом для оркестратора служит движок копирования: запустить перенос между двумя \gls{s3}-бакетами, опросить прогресс, при необходимости отменить.
Всё остальное выполняется через внутренние модули \textit{service-staas}.

\subsubsection{Движок копирования}\label{subsubsec:copying-engine}

Роль движка копирования играет \gls{rclone}~--- свободная утилита для синхронизации данных между \gls{s3}-хранилищами~\cite{rclone-docs}. Оркестратор обращается к нему через управляющий интерфейс.

У такой модели есть два значимых ограничения:
\begin{itemize}
    \item каждый запуск пробегает полную сверку бакета~--- \texttt{rclone sync} не сохраняет состояние между запусками (подраздел~\ref{subsec:manual-migration});
    \item нет механизма блокировок на объектах, через который движок мог бы координироваться с прокси (подраздел~\ref{subsubsec:approach-seamless}).
\end{itemize}

Из-за второго пункта такой движок годится только для варианта с даунтаймом.

\subsection{Архитектура бесшовной миграции}\label{subsec:seamless-arch}

Бесшовная миграция использует ту же общую схему с оркестратором и статусной моделью бакета, но движок копирования заменяется собственным мигратором, а перед бакетом ставится прокси \textit{s3-shifter} (рис.~\ref{fig:seamless-architecture}).

\begin{figure}[H]
    \centering
    \includegraphics[width=0.9\textwidth]{seamless-architecture}
    \caption{Компоненты бесшовной миграции и их взаимодействие}
    \label{fig:seamless-architecture}
\end{figure}

\textit{s3-shifter} включается в путь запроса только на время миграции конкретного бакета и исчезает из неё после её завершения. Пока миграция активна, клиентский трафик идёт через прокси, который в зависимости от типа операции направляет запрос в исходный, целевой или сразу оба бакета; детали маршрутизации рассмотрены далее. Параллельно мигратор фоном переносит объекты, и чтобы исключить гонки, оба компонента согласуются через распределённые блокировки на объектах. \textit{service-staas} управляет жизненным циклом миграции~--- запускает мигратор, инициирует переключение DNS и доводит бакет до конечного статуса.

\subsubsection{Прокси s3-shifter}\label{subsubsec:s3-shifter}

На время миграции \gls{consul}-DNS отдаёт клиенту адрес \textit{s3-shifter} вместо адреса конкретного кластера~--- так весь S3-трафик попадает в прокси.
После завершения миграции DNS-запись возвращается на новый кластер напрямую, и прокси исчезает из пути.

Маршрутизация запросов зависит от типа операции (рис.~\ref{fig:shifter-routing}):
\begin{itemize}
    \item \textbf{PUT}~--- только в целевой бакет; новых данных в исходный больше не приходит, поэтому фоновое копирование сходится;
    \item \textbf{GET}~--- идёт сначала в целевой; при ответе 404 запрос проксируется в исходный;
    \item \textbf{DELETE}~--- в оба бакета (сначала исходный, затем целевой), чтобы исключить переливку объекта обратно мигратором;
    \item \textbf{LIST}~--- листинги обоих бакетов объединяются, при совпадении ключей приоритет имеет целевой;
    \item \textbf{Multipart Upload}~--- \texttt{uploadId} привязан к бакету, поэтому новые загрузки создаются в целевом; операции над MPU, начатым до миграции, при отсутствии \texttt{uploadId} в целевом направляются в исходный~\cite{aws-s3}.
\end{itemize}

\begin{figure}[H]
    \centering
    \includegraphics[width=0.95\textwidth]{shifter-routing}
    \caption{Маршрутизация запросов \textit{s3-shifter} по типам операций}
    \label{fig:shifter-routing}
\end{figure}

Большинство операций (PUT, GET, отдельные части MPU) обрабатывается потоково: тело запроса и ответа сразу пересылается между клиентом и бакетом без накопления в памяти \textit{s3-shifter}.
Буферизация применяется только к запросам с заведомо небольшим телом~--- листингам, метаданным, ответам DELETE.

Чтобы исключить гонку с фоновым копированием, перед каждой модификацией или удалением \textit{s3-shifter} берёт распределённую блокировку на объект и снимает её по завершении операции. Мигратор перед копированием конкретного объекта берёт ту же блокировку (подраздел~\ref{subsubsec:lister-copier}), благодаря чему одновременная запись из двух источников становится невозможной.

Чтения и листинги (GET, LIST) обрабатываются без блокировок: их корректность обеспечивается логикой маршрутизации, а накладные расходы на горячем пути остаются минимальными. Сама блокировка автоматически освобождается при потере владельца, чтобы упавший процесс не оставлял объект заблокированным навсегда.

\subsubsection{Процесс копирования}\label{subsubsec:lister-copier}

Собственный мигратор разделён на две роли~--- Lister и Copier, обмен между которыми идёт через персистентную базу данных с состоянием миграции (рис.~\ref{fig:copy-process}).
Lister отвечает за обнаружение объектов, Copier~--- за их перенос.

\begin{figure}[H]
    \centering
    \includegraphics[width=0.7\textwidth]{copy-process}
    \caption{Поток данных при копировании объектов мигратором}
    \label{fig:copy-process}
\end{figure}

\textbf{Lister} последовательно листает исходный бакет вызовами \texttt{ListObjects}, и для каждого нового ключа создаёт запись в БД со статусом \code{INIT}.
Параллелить эту роль не требуется: листинг S3 уже отдаёт ключи большими страницами, и узким местом обычно становится не он, а копирование.

Согласование между параллельными Copier'ами строится на двух механизмах. Общая БД сериализует выдачу ключей (один ключ в работу выдаётся только одному Copier'у), а распределённые блокировки на объектах страхуют от гонок на уровне самого объекта. Дополнительно к этому каждое назначение ключа фиксирует владельца и срок аренды обработки, а также монотонный \textit{fencing token}: это защищает от классической проблемы распределённых блокировок~\cite{kleppmann-locks}~--- процесс, считающий себя владельцем блокировки, фактически уже мог её потерять из-за паузы (GC, сетевая изоляция).

\textbf{Copier} выбирает из БД очередной ключ со статусом \code{INIT} и переводит его в \code{IN\_PROGRESS}, фиксируя владельца, срок аренды и новый \textit{fencing token}. После этого он берёт распределённую блокировку на объект и, если объекта в целевом бакете нет, потоково читает его из исходного (\texttt{GetObject}) и загружает в целевой через Multipart Upload, разбивая поток на части и отправляя их параллельно.

При финальном обновлении статуса воркер проверяет свой \textit{fencing token}: устаревший воркер после потери аренды не сможет перезаписать состояние, если объект уже перехвачен новым владельцем. Если воркер завершился до финального обновления, по истечении аренды объект возвращается в \code{INIT} и подхватывается другим воркером.

По завершении статус ключа обновляется на одно из терминальных значений и блокировка снимается (рис.~\ref{fig:object-states}):
\begin{itemize}
    \item \code{COPIED}~--- объект скопирован;
    \item \code{EXIST}~--- объект уже присутствовал в целевом бакете (например, был записан клиентом через \textit{s3-shifter}), копирование пропущено;
    \item \code{DELETED}~--- объект удалён после старта миграции и копировать его не требуется.
\end{itemize}

\begin{figure}[H]
    \centering
    \includegraphics[width=0.7\textwidth]{object-states}
    \caption{Статусная модель объекта в БД мигратора}
    \label{fig:object-states}
\end{figure}

Персистентность БД~--- ключевое условие отказоустойчивости: при перезапуске мигратора прогресс не теряется, Lister подхватывает листинг с последней страницы, а Copier продолжает обрабатывать ключи, не помеченные как завершённые.
Та же БД хранит общий статус миграции бакета, который мигратор отдаёт \textit{service-staas}, а тот отражает его в статусной модели платформы (подраздел~\ref{subsubsec:orchestrator}).

Архитектура допускает горизонтальное масштабирование мигратора: общая БД и распределённые блокировки делают одновременную работу нескольких экземпляров безопасной.

\subsection{Эксплуатационные свойства}\label{subsec:cross-cutting}

Сквозные свойства решения, обеспечивающие соответствие нефункциональным требованиям из подраздела~\ref{subsec:nonfunc-req}:
\begin{itemize}
    \item \textbf{идемпотентность.} Все операции мигратора и \textit{s3-shifter} безопасны к повтору: переходы статусов в модели объекта терминальны, незавершённое копирование не оставляет видимого результата в целевом бакете, \texttt{DELETE} по \gls{s3} идемпотентен по протоколу. Это базовое условие для возобновления миграции после сбоя~--- без идемпотентности повторный запуск приводил бы к расхождению данных;
    \item \textbf{восстановление после сбоев.} Прогресс миграции вынесен во внешнее персистентное хранилище, поэтому перезапуск любого компонента продолжает работу с того же места. Удерживаемые блокировки освобождаются автоматически при потере владельца, а механика аренды и \textit{fencing token} исключает потери прогресса при отказе;
    \item \textbf{ограничение нагрузки.} Параллелизм копирования ограничен сверху, чтобы фоновое копирование не насыщало пропускную способность исходного и целевого кластеров. Это удерживает латентность клиентских операций через \textit{s3-shifter} в пределах \gls{slo} платформы (подраздел~\ref{subsec:nonfunc-req}) даже при работе мигратора на полной мощности;
    \item \textbf{наблюдаемость.} Архитектура предусматривает три класса показателей: прогресс миграции (доля скопированных ключей, скорость, ошибки), метрики \gls{slo} клиентских операций через \textit{s3-shifter} и сквозные трассировки для разбора инцидентов на уровне отдельного объекта.
\end{itemize}

\subsection{Выводы по главе}\label{subsec:ch2-summary}

В главе спроектирован механизм миграции \gls{s3}-бакетов в двух режимах: с согласованным даунтаймом и бесшовном.
Базой обоих режимов служит статусная модель бакета в \textit{service-staas} с оркестратором, ведущим миграцию по фазам; для бесшовного режима архитектура дополнена прокси \textit{s3-shifter} и собственным мигратором с ролями Lister/Copier, персистентным хранилищем прогресса и распределёнными блокировками с TTL-арендой и \textit{fencing token}.
Сквозные свойства (идемпотентность, восстановление после сбоев, ограничение нагрузки, наблюдаемость) увязаны с нефункциональными требованиями из подраздела~\ref{subsec:nonfunc-req}.
